% ============================================================
%  Introduction to NLP — Research Project Proposal Template
%  2 pages max (excluding references)
%
%  TOGGLE — change the flag below to switch between versions:
%    \instructionstrue  → full version: grey explanations + blue examples visible
%    \instructionsfalse → clean version: section headers and student text only
% ============================================================
\documentclass[11pt, a4paper]{article}

% --- Version toggle -----------------------------------------
\newif\ifinstructions
\instructionstrue   % <-- change to \instructionsfalse for the clean version

% --- Packages -----------------------------------------------
\usepackage[margin=1in]{geometry}
\usepackage{titlesec}
\usepackage{enumitem}
\usepackage[dvipsnames]{xcolor}
\usepackage{hyperref}
\usepackage{parskip}
\usepackage{booktabs}
\usepackage[authoryear,round]{natbib}

\usepackage{microtype}
\usepackage{fancyhdr}
\usepackage[T1]{fontenc}
\usepackage[utf8]{inputenc}
\usepackage{soul}        % for \hl (yellow highlight)
\usepackage{tabularx}
\usepackage{array}

% --- Colors -------------------------------------------------
\definecolor{exgray}{HTML}{888780}       % explanation text
\definecolor{exblue}{HTML}{185FA5}       % example text
\definecolor{highlight}{HTML}{FFE066}    % yellow highlight for key phrases
\definecolor{sectioncolor}{HTML}{2C2C2A}
\definecolor{rulegray}{HTML}{B4B2A9}

% --- Convenience commands -----------------------------------

% Grey explanation paragraph — hidden in clean version
\newcommand{\explain}[1]{%
  \ifinstructions
    \par\vspace{2pt}%
    {\small\color{exgray} #1}%
    \par\vspace{2pt}%
  \fi
}

% Grey guiding question — hidden in clean version
\newcommand{\guideline}[1]{%
  \ifinstructions
    \par\vspace{3pt}%
    {\small\color{exgray}\textit{#1}}%
    \par\vspace{1pt}%
  \fi
}

% Blue example answer — hidden in clean version
\newcommand{\example}[1]{%
  \ifinstructions
    \par\vspace{1pt}%
    {\small\color{exblue} #1}%
    \par\vspace{4pt}%
  \fi
}

% Yellow-highlighted key phrase — plain text in clean version
\newcommand{\keyphrase}[1]{%
  \ifinstructions
    \sethlcolor{highlight}\hl{#1}%
  \else
    #1%
  \fi
}

% Student fill-in placeholder
\newcommand{\fillhere}[1]{%
  \par\vspace{2pt}%
  \textit{\textcolor{gray}{#1}}%
  \par\vspace{2pt}%
}


% --- Section formatting -------------------------------------
\titleformat{\section}
  {\normalfont\large\bfseries\color{sectioncolor}}
  {\thesection.}{0.5em}{}[\vspace{-8pt}\noindent\rule{\linewidth}{0.2pt}]

% --- Header / Footer ----------------------------------------
\pagestyle{fancy}
\fancyhf{}
\renewcommand{\headrulewidth}{0.4pt}
\fancyhead[L]{\small\textit{TAR --- Research Project Proposal}}
\fancyfoot[C]{\small\thepage}

% ============================================================
\begin{document}

% --- Title block --------------------------------------------
\begin{center}
  {\LARGE\bfseries Working Title}\\[8pt]
\end{center}
\vspace{2pt}

% --- How to use this template (visible in instructor version only) ---
\ifinstructions
\vspace{6pt}
\noindent{\small\color{exgray}
\textbf{How to use this template}

Throughout this template, \textcolor{exgray}{\textbf{grey text}} provides instructions and guiding questions for each section.
\textcolor{exblue}{\textbf{Blue text}} shows concrete examples of what a strong answer looks like --- use it as a reference, do not copy it.
Phrases \keyphrase{highlighted in yellow} are key phrases you can use in your proposal.
Italic grey placeholders such as \textit{\textcolor{gray}{<your text here>}} mark where your own writing should go.

\smallskip
\textbf{Submission}

Write your proposal by filling in the placeholder text throughout the document.
When you are ready to submit, change \texttt{\textbackslash instructionstrue} to \texttt{\textbackslash instructionsfalse} at the top of the \texttt{.tex} file (line 13), recompile, and submit the resulting \textbf{PDF}.
The proposal must not exceed \textbf{3 pages} excluding references.

\smallskip
\textbf{References}

\explain{References are \textbf{not} counted toward the 2-page limit. Add your references to the \texttt{references.bib} file. Each entry has a unique key (e.g.\ \texttt{devlin2019bert}) that you use in \texttt{\textbackslash citep\{\}} and \texttt{\textbackslash citet\{\}} commands. The reference list below is generated automatically from that file when you compile.}

\explain{How to cite in the text:
\begin{itemize}[leftmargin=*, itemsep=1pt, topsep=2pt]
  \item Parenthetical: \texttt{\textbackslash citep\{key\}} produces \textcolor{exblue}{(Author, Year)}
  \item Author as subject: \texttt{\textbackslash citet\{key\}} produces \textcolor{exblue}{Author (Year)}
  \item Multiple references: \texttt{\textbackslash citep\{key1,key2\}} produces \textcolor{exblue}{(Author1, Year1; Author2, Year2)}
\end{itemize}}


\noindent\textcolor{rulegray}{\rule{\linewidth}{0.4pt}}
\vspace{4pt}
\fi

% ============================================================
% SECTION 1 — BASICS
% ============================================================
\section{Basics}

\begin{tabular}{@{}ll}
  \textbf{Project topic:} & \fillhere{<Project topic assigned from the TAR 2026 topic list.>}   \\[6pt]
  \textbf{Team:} & \fillhere{<Team name.>} \\[2pt]
  \textbf{Team members:}  & \fillhere{<Full Name>; <JMBAG>} \\[2pt]
                          & \fillhere{<Full Name>; <JMBAG>} \\[2pt]
                          & \fillhere{<Full Name>; <JMBAG>} \\
\end{tabular}
% ============================================================
% SECTION 2 — ELEVATOR PITCH
% ============================================================
\section{Elevator Pitch}

\explain{Describe your project in 2--3 sentences as if you had 30 seconds to explain it to someone unfamiliar with the topic. Name the \textbf{problem}, the \textbf{approach}, and the expected \textbf{contribution} --- in that order. Write this section \textbf{last}, even though it appears first. It is the abstract of your abstract.}

\guideline{(1) What problem does your work address, and why does it matter?}
\example{``Despite the rapid growth of NLP benchmarks, most evaluate models in English only, leaving multilingual generalization largely unmeasured.''}

\guideline{(2) What is your approach?}
\example{``We investigate whether cross-lingual transfer from pre-trained multilingual models can close this gap on low-resource Slavic languages.''}

\guideline{(3) What will your work contribute?}
\example{``Our study produces a new benchmark and a set of empirical findings that will inform future multilingual model design.''}

\fillhere{
With the advent of social media, the emergence of conspiracy theories has become frequent.
We investigate the phenomenon of conspiracy in the context of Large Language Models.
Our study will produce new empirical findings that help advance this field.}

% ============================================================
% SECTION 3 — MOTIVATION
% ============================================================
\section{Motivation}

\explain{If you can bring references to support your claims, do that. Feel free to use the phrases \keyphrase{highlighted in yellow} as is or their paraphrases.}

\guideline{Start big. What is the broad, general problem you are addressing? The bigger the better.}
\example{``The number of NLP publications has grown exponentially over the past decade \citep{joshi2020state, conneau2020xlmr}, and language diversity remains a fundamental challenge for the field.''}

\guideline{What particular, focused problem are you addressing?}
\example{``Despite this growth, the \keyphrase{main goal addressed by this paper} is understanding to what extent current pre-trained models generalize across typologically different languages without task-specific fine-tuning data.''}

\guideline{Why is this an important problem?}
\example{``\keyphrase{It is an important goal because} the majority of the world's languages are low-resource, and excluding them from NLP research perpetuates a language divide with real societal consequences \citep{joshi2020state, lauscher2020zero}.''}

\guideline{Why is this a non-trivial, hard problem?}
\example{``\keyphrase{This is hard to do because} language typology, script systems, and data availability interact in complex ways that make it difficult to predict when and why transfer succeeds or fails.''}

\fillhere{<Your motivation, following the four-step structure above. Approx.\ 150 words.>}

% ============================================================
% SECTION 4 — GAP
% ============================================================
\section{Gap}

\explain{Show what was done and what is missing. Think wider than just your topic: what is the broad class of problems your work is related to? Usually you will identify 3--4 relevant papers, which you will use as a starting point for your project. \textbf{This step is critical.} Do a thorough and systematic related work analysis \textit{before} starting work on the project, to avoid solving a non-existing problem. The gap must be clear; continue monitoring related work as you proceed. If things change, revise your plan.}

\guideline{What has been done in the past?}
\example{``Pre-trained multilingual models such as mBERT and XLM-R have shown strong cross-lingual transfer on high-resource language pairs \citep{devlin2019bert, conneau2020xlmr, hu2020xtreme}.''}

\guideline{What is the problem with existing work?}
\example{``However, evaluations are dominated by a small set of languages \citep{joshi2020state}. Low-resource Slavic languages are rarely included \citep{hu2020xtreme, lauscher2020zero}. Existing benchmarks do not control for domain shift, making it hard to isolate the contribution of cross-lingual transfer from domain mismatch \citep{conneau2020xlmr}. Thus, no single existing study provides a comprehensive picture of transfer in this setting.''}

\guideline{What is the missing part?}
\example{``Paper X evaluates cross-lingual transfer but only on high-resource languages. Paper Y covers low-resource languages but lacks domain control. Paper Z does both, but the data is not made public. \keyphrase{To the best of our knowledge}, our work is the first to evaluate cross-lingual transfer systematically across seven low-resource Slavic languages under controlled domain conditions.''}

\fillhere{
Analysis of the semantic structure of conspiracy theories has been a common problem in psychoanalysis. \citep{fort2023bigfoot} addresses this problem by training the model on data in which topical words have been masked. The issue with this paper is the usage of outdated models such as BART and SVM. There is potential to build on their work by using LLMs. \citep{platt2022toward} is based on a similar idea. The goal of that paper is to differentiate between conspiracy theory text and more general misinformation text. Another interesting research direction deals with pretraining . \citep{levy2021investigating} investigates GPT- 2's ability to generate conspiracy text. They explain this is a consequence of conspiracy text being deeply rooted in the pretraining data. The problem with this paper is that it deals with the outdated GPT-2 .
}

% ============================================================
% SECTION 5 — RESEARCH QUESTIONS \& CONTRIBUTIONS
% ============================================================
\section{Research Questions \& Contributions}

\explain{List your research questions (RQs). An RQ is a question your study will answer --- what will it add to our knowledge compared to related work? Good RQs start with: ``What is the relationship between\ldots'', ``What are the factors that affect\ldots'', ``What types of X exist\ldots'' (exploratory), ``How to do X\ldots'' (engineering). Avoid ``Can X do Y?'' unless you are confident it can, or will provide strong theoretical or empirical evidence that it cannot.}

\explain{For each RQ, list your planned \textbf{contributions} --- the tangible, reusable results your study will produce. A contribution can be a framework, a dataset, a model, a thorough analysis, or a research agenda. You will typically have 1--3 contributions.}

\example{``To enable the study of cross-lingual transfer in low-resource Slavic languages, this paper addresses three research questions and makes the following contributions:

\textbf{RQ1:} How well does zero-shot cross-lingual transfer from XLM-R generalize to low-resource Slavic languages on NER?\\
\textit{Contribution:} We report a systematic evaluation across seven languages and three domains, providing the most comprehensive picture to date of transfer performance in this setting.
More efficient.

\textbf{RQ2:} What linguistic factors predict transfer success or failure?\\
\textit{Contribution:} We introduce a diagnostic framework that correlates model errors with morphological complexity, script type, and training data size.

\textbf{RQ3:} How to improve transfer for the lowest-resource languages?\\
\textit{Contribution:} We propose and evaluate a lightweight adaptation strategy using automatically translated pivot data. Our results will pave the path for more inclusive multilingual NLP benchmarking.''}

\noindent\textbf{RQ1:} 
\fillhere{
  Is it possible to detect a conspiracy theory if given a text in which
  actors, actions, etc. are masked?}
\textit{Contribution:}
\fillhere{This will give us a better insight in conspiracy theory texts. We will
determine if conspiracy theries follow specific forms that are independent of
factual truth and that can be used to differentiate between critically written
text and a conspiracy theory.}

\noindent\textbf{RQ2:} \fillhere{
Should the pretraining set contain conspiracy theories?
}
\textit{Contribution:} \fillhere{
We aim to determine if additional pretraining on harmful data such as
conspiracy theories can be done in a safe way which also improves the model.
If proven possible this could lead to other research regarding AI safety}

\noindent\textbf{RQ3:} \fillhere{
Can we addionally fine-tune the model to be resistant to sarcasm?
}
\textit{Contribution:} \fillhere{<What tangible result will this produce?>}

\noindent\textbf{RQ4:} \fillhere{
How does learning conspiracy theories affect models performance on other tasks?
}
\textit{Contribution:} \fillhere{<What tangible result will this produce?>}

% ============================================================
% SECTION 6 — EXPERIMENTAL SETUP
% ============================================================
\section{Experimental Setup}

\explain{Briefly describe your experimental plan across the four areas below. Each choice must be \textbf{justified} with respect to your RQs. For example,  saying ``we will use BERT'' without explaining why teaches nothing. This section should be concrete enough to show the pipeline is feasible, but you do not need a full methods section. Close with a short paragraph on risks.}

\paragraph{Datasets.}
\explain{What data will you use? Is it publicly available? Describe its key characteristics (size, domain, language, annotation scheme). If you plan to collect or synthesize data, briefly explain how.}
\example{``We will use the WikiAnn dataset \citep{pan2017wikiann}, which provides NER annotations for 176 languages including all seven target Slavic languages. It is publicly available, large-scale, and has been used as a standard benchmark in prior cross-lingual transfer work \citep{hu2020xtreme, lauscher2020zero}, making our results directly comparable.''}
\fillhere{<Your dataset description and justification.>}

\paragraph{Models \& Baselines.}
\explain{What systems will you build, fine-tune, or prompt? What baselines will you compare against and why? Justify why these models are appropriate for your RQs.}
\example{``Our primary model is XLM-R\textsubscript{large} \citep{conneau2020xlmr}, chosen for its strong multilingual coverage of Slavic scripts. We compare against mBERT \citep{devlin2019bert} as a weaker multilingual baseline and a language-specific monolingual model as an upper-bound oracle. This setup directly addresses RQ1 by isolating the effect of multilingual pre-training scale.''}
\fillhere{<Your model and baseline description and justification.>}

\paragraph{Evaluation.}
\explain{What metrics will you use to measure performance? Why are these metrics appropriate for your RQs? Note any known limitations of these metrics.}
\example{``We report span-level F\textsubscript{1} following the CoNLL evaluation standard, which is the de facto metric for NER and enables direct comparison with prior work. We also report per-language and per-entity-type breakdowns to address RQ2. We note that F\textsubscript{1} does not distinguish between boundary and label errors, which we will supplement with a qualitative error analysis.''}
\fillhere{<Your evaluation metrics and justification.>}

\paragraph{General Methodology.}
\explain{Describe your experimental pipeline in 3--5 steps. Walk the reader through how you will go from raw data to an answer for each RQ. This does not need to be exhaustive --- just enough to show the plan is concrete and well thought through.}
\example{``(1) Preprocess and reformat WikiAnn splits for all seven target languages. (2) Fine-tune XLM-R and mBERT on the English training split only (zero-shot transfer setup). (3) Evaluate all models on each target language test set and record F\textsubscript{1} per language and entity type (RQ1). (4) Correlate error rates with language-level features such as morphological complexity and script type (RQ2). (5) Apply the pivot-data adaptation strategy and re-evaluate on the three lowest-resource languages (RQ3).''}
\fillhere{<Your methodology, described in 3--5 steps.>}

\paragraph{Risks.}
\explain{What could go wrong? For each risk, state the impact on your RQs, and your mitigation or fallback plan.}
\example{``\textit{Data availability:} WikiAnn coverage for two target languages is sparse ($<$1k sentences); if quality is insufficient we will fall back to silver data generated via projection from a related high-resource language. \textit{Compute:} Fine-tuning XLM-R\textsubscript{large} requires significant GPU memory; if unavailable we will use the base variant, which has been shown to close 85\% of the gap at a fraction of the cost.''}
\fillhere{<1--3 concrete risks with mitigations.>}

% ============================================================
% SECTION 7 — JUSTIFICATION FOR API CREDITS
% ============================================================
\section{Justification for API Credits}

\explain{Write your justification for using the API credits as a mini-pitch --- be \textbf{specific}. ``We need the API to generate synthetic training data for low-resource language X because no open dataset exists'' is compelling. ``We want to try a large model'' is not. You must also include a \textbf{fallback plan} --- your project must be viable regardless of whether credits are awarded.}

\guideline{Use case: what specifically would you use the API for, and which RQ does it serve?}
\example{``We plan to use the API to generate silver-standard NER annotations for three languages for which no labeled data exists, directly enabling RQ3.''}

\guideline{Alternative plan: if you do not receive credits, what is your fallback?}
\example{``We will restrict the study to the four languages for which labeled data already exists, narrowing the scope of RQ3 but keeping the core contributions intact.''}

\fillhere{<Your justification, following the structure above. Approx.\ 150 words.>}

% ============================================================
% REFERENCES
% ============================================================

\vspace{4pt}
\bibliographystyle{apalike}
\bibliography{references}

\end{document}
